\section{INTRODUCTION}

\subsection{Background}
In contemporary healthcare institutions, the widespread adoption of Electronic Health Record (EHR) and Electronic Medical Record (EMR) systems has resulted in the continuous generation of massive volumes of clinical narratives. Healthcare practitioners generate extensive unstructured text daily, including admission notes, preoperative consultations, surgical operative summaries, progress notes, and diagnostic evaluations. Accurately categorizing these unstructured clinical notes into specific medical specialties (such as Cardiology, Neurology, Orthopedics, Gastroenterology, Urology, and Ophthalmology) is an indispensable prerequisite for hospital workflow automation, patient file routing, diagnostic database indexing, and clinical decision support.

Traditionally, document indexing relies on manual review by medical coders and health information management personnel. With escalating patient volumes in modern medical centers, manual categorization has become a severe administrative bottleneck, introducing substantial latency, high labor costs, and inter-annotator classification discrepancies. Natural Language Processing (NLP) paired with Machine Learning (ML) provides an automated computational paradigm capable of rapidly extracting semantic signals from clinical narratives, standardizing specialty categorization, and routing medical documentation in real time.

\subsection{Problem Statement}
Manual categorization of medical transcription narratives into appropriate clinical specialties is labor-intensive, time-consuming, and prone to subjective diagnostic interpretation. Unstructured clinical notes present formidable linguistic complexities, including dense domain-specific medical terminology, non-standard abbreviations, variable document lengths, and anatomical lexical overlap across adjacent clinical sub-disciplines.

Furthermore, authentic medical transcription datasets suffer from severe class imbalance and noisy labeling. Datasets such as the Kaggle Medical Transcriptions (MTSamples) corpus initially group notes into broad procedural or administrative categories (such as General Surgery, SOAP Notes, and Consultative Summaries), which overlap heavily with organ-specific departments. A critical limitation of conventional machine learning systems is their closed-set assumption: they assume that every incoming clinical record belongs to one of the pre-defined target classes. In real-world hospital environments, out-of-scope clinical narratives (such as Dermatology, Psychiatry, or Pediatrics) or non-medical text are routinely encountered. Closed-set classifiers mistakenly force-fit these out-of-domain cases into an arbitrary specialty, compromising clinical integrity. Additionally, deep learning transformer models (such as BERT or BioBERT) incur excessive GPU computational overhead, rendering them impractical for lightweight, instantaneous web deployment.

\subsection{Need for the System}
Automating medical specialty classification through efficient statistical natural language processing and ensemble machine learning addresses critical clinical needs:
\begin{itemize}
    \item \textit{Streamlined Clinical Triage}: Eliminates manual document sorting backlogs, enabling automated indexing and immediate department routing upon clinical dictation.
    \item \textit{Rapid Emergency Retrieval}: Empowers attending physicians to quickly access historical specialty-specific operative notes during acute patient emergencies.
    \item \textit{Open-Set Patient Safety}: Implements a confidence-gated rejection mechanism that identifies out-of-scope or ambiguous medical documents, safely routing them to manual clinical review rather than producing hazardous misclassifications.
    \item \textit{Explainable NLP Representation}: Provides algorithmic transparency through Term Frequency--Inverse Document Frequency (TF-IDF) feature attribution, allowing clinicians to inspect the exact keywords driving model decisions.
    \item \textit{Computationally Efficient Web Deployment}: Delivers instantaneous inference on standard CPU architectures through lightweight classical estimators, hosted via an accessible web platform with direct PDF report ingestion.
\end{itemize}

\subsection{Objectives}
The primary objective of this project is to design, develop, and deploy a robust, end-to-end automated Medical Specialty Classification system (\textbf{MedSpecialty AI Platform}) using sublinear TF-IDF vectorization and an Ensemble Machine Learning architecture with open-set rejection capabilities. The specific objectives are:
\begin{enumerate}
    \item To conduct comprehensive Exploratory Data Analysis (EDA) on the raw MTSamples medical transcriptions corpus (4,999 records across 40 categories) and identify noisy, overlapping, and administrative document classes.
    \item To curate a clinically sound, mutually exclusive 8-specialty benchmark dataset comprising 1,663 records across core hospital departments.
    \item To build a modular clinical NLP preprocessing pipeline incorporating lowercasing, medical punctuation handling, tokenization, stopword removal, and morphological lemmatization.
    \item To extract discriminative unigram and bigram features using an 8,000-dimensional sublinear TF-IDF vectorizer with $L_2$ Euclidean normalization.
    \item To train and evaluate candidate supervised classifiers with distinct inductive biases: Linear Support Vector Machine (Linear SVM with Platt probability calibration), Logistic Regression, Random Forest, and Multinomial Naive Bayes.
    \item To construct a Soft Voting Ensemble Classifier utilizing calibrated weights [2:2:1:1] to synthesize probabilistic consensus across candidate models.
    \item To formulate an Open-Set Confidence Rejection mechanism with a calibrated safety threshold ($\tau = 0.48$) establishing a 9th operational class (\textit{``Other / Requires Review''}) for out-of-scope clinical records.
    \item To develop and deploy an interactive, production-grade web application using Flask, supporting drag-and-drop PDF text extraction via PyMuPDF, real-time transcription editing, and 7-phase dataflow visualization.
\end{enumerate}

\subsection{Scope}
The scope of this project encompasses:
\begin{itemize}
    \item \textit{Dataset Scope}: Evaluates 4,999 raw clinical transcription records from the Kaggle MTSamples corpus, filtering out non-specialty administrative document formats (SOAP notes, Consults, Discharge Summaries) and overlapping umbrellas (General Surgery) to establish a curated benchmark of 1,663 records across eight core hospital departments:
    \begin{enumerate}
        \item Cardiovascular / Pulmonary
        \item Orthopedic
        \item Gastroenterology
        \item Neurology
        \item Urology
        \item Obstetrics / Gynecology
        \item ENT -- Otolaryngology
        \item Ophthalmology
    \end{enumerate}
    \item \textit{Feature Engineering Scope}: Focuses on unigram and bigram TF-IDF representations bounded at 8,000 features, applying logarithmic term frequency dampening ($1 + \log(\text{TF})$) to prevent high-frequency term domination.
    \item \textit{Algorithmic Scope}: Investigates four classical machine learning models (Linear SVM, Logistic Regression, Random Forest, and Multinomial Naive Bayes), a Soft Voting Ensemble, and a threshold-based open-set rejection gate.
    \item \textit{Deployment Scope}: Delivers a fully responsive web application (\textbf{MedSpecialty AI Platform}) built with Flask and Vanilla CSS/JavaScript, featuring in-memory PDF parsing (PyMuPDF) and full pipeline explainability without requiring external GPU infrastructure.
\end{itemize}

\subsection{Project Overview}
The project follows a structured engineering workflow:
\begin{enumerate}
    \item \textit{Data Ingestion \& Exploratory Profiling}: Loading raw MTSamples CSV data, analyzing missing attributes, document length distributions, and category imbalances.
    \item \textit{Dataset Curation \& Cleaning}: Pruning null records, resolving category ambiguities, and extracting enriched clinical text (combining consultation abstracts and transcription narratives).
    \item \textit{NLP Preprocessing}: Standardizing text via regex cleaning, lowercasing, punctuation stripping, tokenization, stopword removal, and WordNet morphological lemmatization.
    \item \textit{Feature Extraction}: Vectorizing preprocessed narratives into an 8,000-dimensional sparse feature space using sublinear TF-IDF scaling.
    \item \textit{Model Development \& Calibration}: Training candidate estimators on stratified training partitions (80\% train / 20\% test) and applying Platt scaling via \texttt{CalibratedClassifierCV} to Linear SVM.
    \item \textit{Ensemble Integration}: Aggregating model posteriors using weighted soft voting ($2 \times \text{SVM} + 2 \times \text{LR} + 1 \times \text{RF} + 1 \times \text{MNB}) / 6$.
    \item \textit{Open-Set Confidence Gating}: Applying decision threshold $\tau = 0.48$ to classify cases into confirmed specialties or route them to \textit{``Other / Requires Review''}.
    \item \textit{Web Platform Deployment}: Serializing artifacts with Joblib and serving predictions via the MedSpecialty AI Platform with PDF ingestion.
\end{enumerate}

\subsection{Expected Outcome}
The project delivers the following primary outcomes:
\begin{itemize}
    \item A curated, high-quality benchmark dataset of 1,663 medical transcriptions across eight clinical specialties.
    \item An empirical comparative analysis confirming that Linear SVM achieves the highest individual accuracy (88.59\%, Macro F1: 0.8971) and that the Soft Voting Ensemble achieves superior probability calibration (87.99\%, Macro F1: 0.8898).
    \item A verified open-set rejection mechanism that accurately routes unsupported clinical records (e.g., Dermatology, Psychiatry) to the \textit{``Other''} category.
    \item A fully functional, production-ready web application providing real-time PDF extraction and clinical NLP explainability.
    \item Comprehensive academic documentation adhering to MCA departmental standards.
\end{itemize}
